\documentclass[12pt]{article}

\usepackage[margin=1in]{geometry}
\usepackage{longtable}

\title{\textbf{Intelligent HR Analysis System}\
\large Software Design Document}
\author{Suijian Li \ Freddy Montano \ Anthony Perez \ Gary Wu \ Hein Thant}
\date{April 2026}

\begin{document}

%---------------- TITLE PAGE ----------------%
\begin{titlepage}
\centering
\vspace*{2cm}

{\LARGE \textbf{Intelligent HR Analysis System} \par}
\vspace{1cm}

{\Large Software Design Document \par}
\vspace{2cm}

{\large
Suijian Li \
Freddy Montano \
Anthony Perez \
Gary Wu \
Hein Thant
\par}

\vfill

{\large April 2026 \par}

\end{titlepage}

%---------------- TABLE OF CONTENTS ----------------%
\newpage
\tableofcontents
\newpage

%---------------- VERSION HISTORY ----------------%
\section*{Version History}
\addcontentsline{toc}{section}{Version History}

\begin{longtable}{|l|l|l|l|}
\hline
\textbf{User} & \textbf{Date} & \textbf{Reason for Changes} & \textbf{Version} \
\hline
Hein T, Suijian L & [5/8/2026] & Snapshot 1 & 1.0 \
\hline
Freddy M & [5/9/2026] & Snapshot 2 & 2.0 \
\hline
Anthony P & [5/10/2026] & Snapshot 3 & 3.0 \
\hline
Gary W & [5/11/2026] & Snapshot 4 & 4.0 \
\hline
\end{longtable}

\newpage

%---------------- 1 INTRODUCTION ----------------%
\section{Introduction}

\subsection{Purpose}
The purpose of this document is to provide a detailed description of the Intelligent HR Analysis System. This document outlines the system architecture, user interface, and overall system design.

\subsection{Scope}
This document covers the design of the Intelligent HR Analysis System. It serves as a reference for understanding how the system is structured and how components interact.

\subsection{Intended Audience}
\begin{itemize}
\item Development team
\item Project managers
\item Stakeholders
\end{itemize}

\subsection{Overview}
The Intelligent HR Analysis System is a web-based platform designed to automate HR workflows, track employee performance, and provide analytics dashboards for decision-making.

\subsection{References}
\begin{itemize}
\item Project documentation
\item Course materials
\end{itemize}

\subsection{Definitions, Acronyms, and Abbreviations}
See Glossary section.

%---------------- 2 SYSTEM ARCHITECTURE ----------------%
\section{System Architecture}

\subsection{Workflow}
The workflow of the system consists of the following steps:
\begin{itemize}
\item HR inputs employee data into the system
\item The system processes employee performance reviews
\item Supervisors review and update employee records
\item Data is stored in the system database
\item The system generates analytics and reports
\end{itemize}

\subsection{Site Breakdown}
The system consists of the following components:
\begin{itemize}
\item Web Application – Interface for HR and supervisors
\item Backend System – Handles logic and workflows
\item Database – Stores employee data
\item Analytics Module – Displays reports and insights
\item Notification System – Sends alerts and reminders
\end{itemize}

\subsection{Architecture Overview}
The system follows a layered architecture:
\begin{itemize}
\item \textbf{View (Frontend):} User interface for interaction
\item \textbf{Controller (Backend):} Processes requests and logic
\item \textbf{Model (Database):} Stores and manages data
\end{itemize}

This structure ensures modularity and scalability.

\subsection{Data Flow – Employee Review}
\begin{enumerate}
\item HR enters employee data into the system
\item The system stores data in the database
\item Supervisors access and update performance reviews
\item The system processes updates and stores results
\item Analytics module retrieves data and displays dashboards
\end{enumerate}

%---------------- 3 USER INTERFACE ----------------%
\section{User Interface}

\subsection{How to Use}

\subsubsection{Web Application}
Users access the system through a web browser.

Features include:
\begin{itemize}
\item Dashboard for HR analytics
\item Employee performance tracking
\item Reports and data visualization
\end{itemize}

\subsection{Database Explanation}
The database stores:
\begin{itemize}
\item Employee information
\item Performance reviews
\item Workflow status
\end{itemize}

Frontend retrieves data via REST API requests. Backend stores data using database schemas and integrates with UI components.

The database system may use a relational database such as MySQL or PostgreSQL.

%---------------- LEGAL ----------------%
\section{Legal and Ethical Considerations}
Utilizing existing hosting sites would cover most legal considerations. Ethical considerations primarily depend on the system users.

%---------------- GLOSSARY ----------------%
\section*{Glossary}
\addcontentsline{toc}{section}{Glossary}

HR -- Human Resources \
UI -- User Interface \
DB -- Database

%---------------- REFERENCES ----------------%
\section*{References}
\addcontentsline{toc}{section}{References}

\begin{itemize}
\item \texttt{https://ascent.cysun.org/project/project/view/251}
\item Course materials and project guidelines
\end{itemize}

\end{document}
